\section{Which losses are included?}
\subsection{Conduction loss}
An enhanced GaN channel behaves approximately as a resistance. Its instantaneous
loss is $i^2R_{DS(on)}$, so RMS current must be used when averaging. In a
balanced three-phase inverter with $N$ ideally sharing devices per position,
the total is $3I_{phase,rms}^2R_{DS(on)}(T)/N$, with ripple added in quadrature.
The device is warmer in operation than the 25 C datasheet test, and the model
must account for this resistance increase. Unequal sharing and dynamic
on-resistance can increase actual loss beyond this estimate.

\subsection{Switching loss: overlap and output capacitance}
During turn-on and turn-off, nonzero drain voltage and current coexist.
Their overlap produces dissipation. A triangular estimate is
$E_{overlap}\simeq V_{dc}|I|(t_{on}+t_{off})/2$. The assumed overlap intervals
are not the gate-driver command rise/fall time or propagation delay.
The EPC2361 datasheet does not tabulate $E_{on}$ and $E_{off}$ for the intended
current, voltage, resistance, temperature and layout; consequently the initial
analytical calculation uses assumed 5 ns intervals rather than measured energies.

Output capacitance is nonlinear. The stored energy is
\[E_{oss}(V)=\int_0^V v C_{oss}(v)\,dv,\]
not generally $C_{oss}(V)V^2/2$ using the single-point capacitance.
Use the energy curve in Figure 6, or the energy-equivalent capacitance at its
specified voltage. The analytical complementary-device fraction is an adjustable
energy allowance, not a proven commutation-loss partition. Circuit conditions
determine how much stored energy is recovered or dissipated.
Once calibrated switching energies replace this approximation, remove the
overlap/Coss terms already contained in those energies.

\subsection{Deadtime: reverse channel conduction}
During deadtime both gate commands are off, but inductive current must continue.
An enhancement-mode GaN device conducts in its third quadrant, with a reverse
voltage drop that depends on current, temperature and off-state gate voltage.
This plays the freewheeling role often associated with a MOSFET body diode,
but EPC2361 has no silicon-style PN body diode and its datasheet lists zero
reverse-recovery charge. Zero $Q_{rr}$ does not mean zero deadtime or capacitive loss.
For two deadtime intervals per leg per carrier cycle,
$P_{dead}\simeq 6f_s t_d V_{SD}\overline{|I_{phase}|}$ across three legs.
The 2 V reverse drop is an explicit first estimate; it must be replaced with a
reading of Figure 8 at the actual per-device current, or a validated simulation.

\subsection{Gate drive and remaining module losses}
The gate supply delivers approximately $Q_gV_g$ per switching cycle per device,
giving $6NQ_gV_gf_s$. This is dissipated across driver/gate resistances, and is
not all heat in the transistor junction. The model counts it in efficiency but
does not apply it to the transistor cold plate. PCB/copper, capacitors, driver
quiescent power, control and auxiliaries require a separate loss allocation.

\section{Where the parameters come from}
The supplied datasheet is revision 2.4, 20 July 2026. Values below explain the
default configuration; the input JSON and saved summary record the actual run.
\begin{center}\small
\begin{tabular}{p{0.21\linewidth}p{0.32\linewidth}p{0.37\linewidth}}
Parameter & Datasheet source & Treatment in the model \\
\hline
$R_{DS(on)}$ & Page 1: 0.75 mOhm typical, 1 mOhm maximum at 25 C, 5 V gate, 50 A &
Use 1 mOhm and about 1.65 temperature factor at 125 C from Figure 9. Factor is a typical visual reading. \\
$Q_g$, $R_g$ & Page 2: 28 nC typical, 34 nC maximum at 50 V; internal $R_g=0.4$ ohm &
Analytical gate power uses 34 nC. SPICE includes internal gate resistance; external and driver resistance are added once. \\
$E_{oss}$ & Page 3, Figure 6; page 2 $C_{oss,ER}=1419$ pF over 0--50 V &
About 3.2 microjoule visually at 75 V; 50 V value is about 1.77 microjoule. Do not scale constant capacitance blindly to other voltages. \\
$V_{SD}$ & Page 1: 1.6 V at only 0.5 A; Figure 8 provides reverse-current curves &
The 0.5 A point is not a high-current drop. Default 2 V is an assumption pending calibration. \\
$E_{on}$, $E_{off}$ & No matching energy tables; page 6 switching traces use a particular driver and layout &
Initial overlap times are assumed; the separate double-pulse sweep extracts explicitly defined model energies. \\
$R_{jc}$, temperature & Page 1: 0.2 K/W typical to top case, 150 C absolute limit &
Use your 125 C design ceiling; TIM and cold plate are separate assumed resistances. \\
Parasitics and driver & Not specified for the user's PCB &
Sweep loop inductance and external gate resistance. Replace driver impedance and layout assumptions with real data. \\
\end{tabular}
\end{center}
\clearpage
